% ============================================================
%  ECON306: W11D1 — Signalling Theory (Spence 1973);
%          Separating Equilibria
%  Date: Wednesday 30 September 2026
%  SOURCE: old w13_week13 (signalling section) from FULL deck
%    (lines 3856–3925) --- first lecture portion.
% ============================================================

\documentclass[aspectratio=169, 12pt]{beamer}

% ============================================================
%  ECON306: Shared Preamble — The Economics of Health and Education
%  University of Otago — Semester 2, 2026
%
%  SHARED BY ALL ECON306 DAY DECKS.
%  Input this file immediately after \documentclass in each deck:
%    % ============================================================
%  ECON306: Shared Preamble — The Economics of Health and Education
%  University of Otago — Semester 2, 2026
%
%  SHARED BY ALL ECON306 DAY DECKS.
%  Input this file immediately after \documentclass in each deck:
%    % ============================================================
%  ECON306: Shared Preamble — The Economics of Health and Education
%  University of Otago — Semester 2, 2026
%
%  SHARED BY ALL ECON306 DAY DECKS.
%  Input this file immediately after \documentclass in each deck:
%    \input{../Preambles/econ306_preamble.tex}
%  Do NOT include \documentclass or per-deck metadata here.
% ============================================================

% ─── Theme & Colours ────────────────────────────────────────
\usetheme{default}
\usecolortheme{default}
\usefonttheme{professionalfonts}
\definecolor{OtagoBlue}{HTML}{003057}
\definecolor{OtagoGold}{HTML}{A8923A}
\definecolor{TextGrey}{HTML}{3A3A3A}
\definecolor{AccentBlue}{HTML}{2B6CB0}
\definecolor{LightBg}{HTML}{F7F9FC}

\setbeamercolor{frametitle}{fg=OtagoBlue, bg=white}
\setbeamercolor{title}{fg=OtagoBlue}
\setbeamercolor{subtitle}{fg=OtagoGold}
\setbeamercolor{author}{fg=TextGrey}
\setbeamercolor{date}{fg=TextGrey}
\setbeamercolor{institute}{fg=TextGrey}
\setbeamercolor{normal text}{fg=TextGrey, bg=white}
\setbeamercolor{item}{fg=OtagoBlue}
\setbeamercolor{subitem}{fg=AccentBlue}
\setbeamercolor{block title}{bg=OtagoBlue!12, fg=OtagoBlue}
\setbeamercolor{block body}{bg=OtagoBlue!5, fg=TextGrey}
\setbeamercolor{alertblock title}{bg=OtagoGold!20, fg=OtagoBlue}
\setbeamercolor{alertblock body}{bg=OtagoGold!8, fg=TextGrey}
\setbeamercolor{alerted text}{fg=OtagoGold!80!black}
\setbeamercolor{structure}{fg=OtagoBlue}

\setbeamertemplate{frametitle}{%
	\vspace{0.35em}%
	\textbf{\insertframetitle}\par%
	\vspace{0.05em}%
	\textcolor{OtagoGold}{\rule{\textwidth}{0.6pt}}}

\setbeamertemplate{navigation symbols}{}
\setbeamertemplate{footline}{%
	\leavevmode\hbox{%
		\begin{beamercolorbox}[wd=.40\paperwidth,ht=2.5ex,dp=1.5ex,leftskip=.3cm]{author in head/foot}%
			\textcolor{TextGrey}{\footnotesize ECON306 \textbullet\ University of Otago}%
		\end{beamercolorbox}%
		\begin{beamercolorbox}[wd=.47\paperwidth,ht=2.5ex,dp=1.5ex,center]{title in head/foot}%
			\textcolor{TextGrey}{\footnotesize \insertshorttitle}%
		\end{beamercolorbox}%
		\begin{beamercolorbox}[wd=.13\paperwidth,ht=2.5ex,dp=1.5ex,rightskip=.3cm,right]{date in head/foot}%
			\textcolor{TextGrey}{\footnotesize \insertframenumber\,/\,\inserttotalframenumber}%
	\end{beamercolorbox}}\vskip0pt}

\setbeamertemplate{itemize item}{\textcolor{OtagoGold}{\small$\blacktriangleright$}}
\setbeamertemplate{itemize subitem}{\textcolor{AccentBlue}{\small$\triangleright$}}
\setbeamertemplate{enumerate item}{\textcolor{OtagoBlue}{\small\bfseries\insertenumlabel.}}

\setbeamertemplate{title page}{%
	\begin{tikzpicture}[remember picture,overlay]
		\fill[OtagoBlue] (current page.north west)
		rectangle ([yshift=-0.38\paperheight]current page.north east);
	\end{tikzpicture}
	\vspace{0.06\paperheight}
	\begin{beamercolorbox}[wd=\textwidth, sep=0pt]{title}
		{\color{white}\Large\textbf{\inserttitle}}\\[0.4em]
		{\color{OtagoGold}\large\insertsubtitle}
	\end{beamercolorbox}
	\vspace{1.4em}
	\begin{beamercolorbox}[wd=\textwidth, sep=0pt]{author}
		{\color{TextGrey}\normalsize\insertauthor}\\[0.2em]
		{\color{TextGrey}\small\insertinstitute}\\[0.2em]
		{\color{TextGrey}\small\insertdate}
\end{beamercolorbox}}

% ─── Packages ──────────────────────────────────────────────
\usepackage{tikz}
\usepackage{booktabs}
\usepackage{array}
\usepackage{graphicx}
\usepackage{hyperref}
\usepackage{amsmath}
\usepackage{amssymb}
\usepackage{colortbl}
\usepackage{multirow}
\usepackage{xcolor}
\hypersetup{colorlinks=true, linkcolor=AccentBlue, urlcolor=AccentBlue,
	citecolor=AccentBlue, pdfauthor={Austin Henderson},
	pdftitle={ECON306 S2 2026}}

% ─── Images folder ─────────────────────────────────────────
%  Place all slide images in the 'images/' subfolder next to each deck.
%  Usage: \includegraphics[width=0.85\textwidth]{figure_name}
\graphicspath{{images/}}

% ─── Reusable macros ───────────────────────────────────────
% Recap slide at start of each lecture
\newcommand{\recapslide}[1]{%
	\begin{frame}{Recap: What Did We Cover Last Time?}
		\begin{block}{Key points from last lecture}
			\begin{itemize}#1\end{itemize}
		\end{block}
\end{frame}}

% Plan slide at start of each lecture
\newcommand{\planslide}[2]{%
	\begin{frame}{Today's Plan — #1}
		\begin{columns}
			\column{0.6\textwidth}
			\begin{itemize}#2\end{itemize}
			\column{0.38\textwidth}
			\begin{alertblock}{Reading}
				\footnotesize See Brightspace (Aoroa) for required readings before this lecture.
			\end{alertblock}
		\end{columns}
\end{frame}}

% Summary slide at end of each lecture
\newcommand{\summaryside}[1]{%
	\begin{frame}{Lecture Summary}
		\begin{exampleblock}{Key takeaways}
			\begin{itemize}#1\end{itemize}
		\end{exampleblock}
		\vspace{0.5em}
		{\footnotesize\textcolor{OtagoGold}{\textbf{Brightspace (Aoroa):}} slides, readings, and any announcements posted there.}
\end{frame}}

% NZ spotlight box
\newcommand{\nzbox}[2]{%
	\begin{alertblock}{\includegraphics[height=0.8em]{nz_flag_icon}\hspace{0.3em} NZ Spotlight: #1}
		#2
\end{alertblock}}

% Tutorial callout
\newcommand{\tutorialbox}[1]{%
	\begin{block}{\textbf{Tutorial This Week}}
		#1
\end{block}}

% ─── Section dividers ──────────────────────────────────────
\AtBeginSection[]{
	\begin{frame}[plain]
		\begin{tikzpicture}[remember picture,overlay]
			\fill[OtagoBlue!95] (current page.north west)
			rectangle (current page.south east);
			\node[anchor=center, text=white, font=\Large\bfseries,
			text width=0.9\paperwidth, align=center]
			at (current page.center) {\insertsectionhead};
			\node[anchor=south, text=OtagoGold, font=\normalsize,
			text width=0.9\paperwidth, align=center]
			at ([yshift=2.5cm]current page.south) {ECON306 — University of Otago};
		\end{tikzpicture}
\end{frame}}

%  Do NOT include \documentclass or per-deck metadata here.
% ============================================================

% ─── Theme & Colours ────────────────────────────────────────
\usetheme{default}
\usecolortheme{default}
\usefonttheme{professionalfonts}
\definecolor{OtagoBlue}{HTML}{003057}
\definecolor{OtagoGold}{HTML}{A8923A}
\definecolor{TextGrey}{HTML}{3A3A3A}
\definecolor{AccentBlue}{HTML}{2B6CB0}
\definecolor{LightBg}{HTML}{F7F9FC}

\setbeamercolor{frametitle}{fg=OtagoBlue, bg=white}
\setbeamercolor{title}{fg=OtagoBlue}
\setbeamercolor{subtitle}{fg=OtagoGold}
\setbeamercolor{author}{fg=TextGrey}
\setbeamercolor{date}{fg=TextGrey}
\setbeamercolor{institute}{fg=TextGrey}
\setbeamercolor{normal text}{fg=TextGrey, bg=white}
\setbeamercolor{item}{fg=OtagoBlue}
\setbeamercolor{subitem}{fg=AccentBlue}
\setbeamercolor{block title}{bg=OtagoBlue!12, fg=OtagoBlue}
\setbeamercolor{block body}{bg=OtagoBlue!5, fg=TextGrey}
\setbeamercolor{alertblock title}{bg=OtagoGold!20, fg=OtagoBlue}
\setbeamercolor{alertblock body}{bg=OtagoGold!8, fg=TextGrey}
\setbeamercolor{alerted text}{fg=OtagoGold!80!black}
\setbeamercolor{structure}{fg=OtagoBlue}

\setbeamertemplate{frametitle}{%
	\vspace{0.35em}%
	\textbf{\insertframetitle}\par%
	\vspace{0.05em}%
	\textcolor{OtagoGold}{\rule{\textwidth}{0.6pt}}}

\setbeamertemplate{navigation symbols}{}
\setbeamertemplate{footline}{%
	\leavevmode\hbox{%
		\begin{beamercolorbox}[wd=.40\paperwidth,ht=2.5ex,dp=1.5ex,leftskip=.3cm]{author in head/foot}%
			\textcolor{TextGrey}{\footnotesize ECON306 \textbullet\ University of Otago}%
		\end{beamercolorbox}%
		\begin{beamercolorbox}[wd=.47\paperwidth,ht=2.5ex,dp=1.5ex,center]{title in head/foot}%
			\textcolor{TextGrey}{\footnotesize \insertshorttitle}%
		\end{beamercolorbox}%
		\begin{beamercolorbox}[wd=.13\paperwidth,ht=2.5ex,dp=1.5ex,rightskip=.3cm,right]{date in head/foot}%
			\textcolor{TextGrey}{\footnotesize \insertframenumber\,/\,\inserttotalframenumber}%
	\end{beamercolorbox}}\vskip0pt}

\setbeamertemplate{itemize item}{\textcolor{OtagoGold}{\small$\blacktriangleright$}}
\setbeamertemplate{itemize subitem}{\textcolor{AccentBlue}{\small$\triangleright$}}
\setbeamertemplate{enumerate item}{\textcolor{OtagoBlue}{\small\bfseries\insertenumlabel.}}

\setbeamertemplate{title page}{%
	\begin{tikzpicture}[remember picture,overlay]
		\fill[OtagoBlue] (current page.north west)
		rectangle ([yshift=-0.38\paperheight]current page.north east);
	\end{tikzpicture}
	\vspace{0.06\paperheight}
	\begin{beamercolorbox}[wd=\textwidth, sep=0pt]{title}
		{\color{white}\Large\textbf{\inserttitle}}\\[0.4em]
		{\color{OtagoGold}\large\insertsubtitle}
	\end{beamercolorbox}
	\vspace{1.4em}
	\begin{beamercolorbox}[wd=\textwidth, sep=0pt]{author}
		{\color{TextGrey}\normalsize\insertauthor}\\[0.2em]
		{\color{TextGrey}\small\insertinstitute}\\[0.2em]
		{\color{TextGrey}\small\insertdate}
\end{beamercolorbox}}

% ─── Packages ──────────────────────────────────────────────
\usepackage{tikz}
\usepackage{booktabs}
\usepackage{array}
\usepackage{graphicx}
\usepackage{hyperref}
\usepackage{amsmath}
\usepackage{amssymb}
\usepackage{colortbl}
\usepackage{multirow}
\usepackage{xcolor}
\hypersetup{colorlinks=true, linkcolor=AccentBlue, urlcolor=AccentBlue,
	citecolor=AccentBlue, pdfauthor={Austin Henderson},
	pdftitle={ECON306 S2 2026}}

% ─── Images folder ─────────────────────────────────────────
%  Place all slide images in the 'images/' subfolder next to each deck.
%  Usage: \includegraphics[width=0.85\textwidth]{figure_name}
\graphicspath{{images/}}

% ─── Reusable macros ───────────────────────────────────────
% Recap slide at start of each lecture
\newcommand{\recapslide}[1]{%
	\begin{frame}{Recap: What Did We Cover Last Time?}
		\begin{block}{Key points from last lecture}
			\begin{itemize}#1\end{itemize}
		\end{block}
\end{frame}}

% Plan slide at start of each lecture
\newcommand{\planslide}[2]{%
	\begin{frame}{Today's Plan — #1}
		\begin{columns}
			\column{0.6\textwidth}
			\begin{itemize}#2\end{itemize}
			\column{0.38\textwidth}
			\begin{alertblock}{Reading}
				\footnotesize See Brightspace (Aoroa) for required readings before this lecture.
			\end{alertblock}
		\end{columns}
\end{frame}}

% Summary slide at end of each lecture
\newcommand{\summaryside}[1]{%
	\begin{frame}{Lecture Summary}
		\begin{exampleblock}{Key takeaways}
			\begin{itemize}#1\end{itemize}
		\end{exampleblock}
		\vspace{0.5em}
		{\footnotesize\textcolor{OtagoGold}{\textbf{Brightspace (Aoroa):}} slides, readings, and any announcements posted there.}
\end{frame}}

% NZ spotlight box
\newcommand{\nzbox}[2]{%
	\begin{alertblock}{\includegraphics[height=0.8em]{nz_flag_icon}\hspace{0.3em} NZ Spotlight: #1}
		#2
\end{alertblock}}

% Tutorial callout
\newcommand{\tutorialbox}[1]{%
	\begin{block}{\textbf{Tutorial This Week}}
		#1
\end{block}}

% ─── Section dividers ──────────────────────────────────────
\AtBeginSection[]{
	\begin{frame}[plain]
		\begin{tikzpicture}[remember picture,overlay]
			\fill[OtagoBlue!95] (current page.north west)
			rectangle (current page.south east);
			\node[anchor=center, text=white, font=\Large\bfseries,
			text width=0.9\paperwidth, align=center]
			at (current page.center) {\insertsectionhead};
			\node[anchor=south, text=OtagoGold, font=\normalsize,
			text width=0.9\paperwidth, align=center]
			at ([yshift=2.5cm]current page.south) {ECON306 — University of Otago};
		\end{tikzpicture}
\end{frame}}

%  Do NOT include \documentclass or per-deck metadata here.
% ============================================================

% ─── Theme & Colours ────────────────────────────────────────
\usetheme{default}
\usecolortheme{default}
\usefonttheme{professionalfonts}
\definecolor{OtagoBlue}{HTML}{003057}
\definecolor{OtagoGold}{HTML}{A8923A}
\definecolor{TextGrey}{HTML}{3A3A3A}
\definecolor{AccentBlue}{HTML}{2B6CB0}
\definecolor{LightBg}{HTML}{F7F9FC}

\setbeamercolor{frametitle}{fg=OtagoBlue, bg=white}
\setbeamercolor{title}{fg=OtagoBlue}
\setbeamercolor{subtitle}{fg=OtagoGold}
\setbeamercolor{author}{fg=TextGrey}
\setbeamercolor{date}{fg=TextGrey}
\setbeamercolor{institute}{fg=TextGrey}
\setbeamercolor{normal text}{fg=TextGrey, bg=white}
\setbeamercolor{item}{fg=OtagoBlue}
\setbeamercolor{subitem}{fg=AccentBlue}
\setbeamercolor{block title}{bg=OtagoBlue!12, fg=OtagoBlue}
\setbeamercolor{block body}{bg=OtagoBlue!5, fg=TextGrey}
\setbeamercolor{alertblock title}{bg=OtagoGold!20, fg=OtagoBlue}
\setbeamercolor{alertblock body}{bg=OtagoGold!8, fg=TextGrey}
\setbeamercolor{alerted text}{fg=OtagoGold!80!black}
\setbeamercolor{structure}{fg=OtagoBlue}

\setbeamertemplate{frametitle}{%
	\vspace{0.35em}%
	\textbf{\insertframetitle}\par%
	\vspace{0.05em}%
	\textcolor{OtagoGold}{\rule{\textwidth}{0.6pt}}}

\setbeamertemplate{navigation symbols}{}
\setbeamertemplate{footline}{%
	\leavevmode\hbox{%
		\begin{beamercolorbox}[wd=.40\paperwidth,ht=2.5ex,dp=1.5ex,leftskip=.3cm]{author in head/foot}%
			\textcolor{TextGrey}{\footnotesize ECON306 \textbullet\ University of Otago}%
		\end{beamercolorbox}%
		\begin{beamercolorbox}[wd=.47\paperwidth,ht=2.5ex,dp=1.5ex,center]{title in head/foot}%
			\textcolor{TextGrey}{\footnotesize \insertshorttitle}%
		\end{beamercolorbox}%
		\begin{beamercolorbox}[wd=.13\paperwidth,ht=2.5ex,dp=1.5ex,rightskip=.3cm,right]{date in head/foot}%
			\textcolor{TextGrey}{\footnotesize \insertframenumber\,/\,\inserttotalframenumber}%
	\end{beamercolorbox}}\vskip0pt}

\setbeamertemplate{itemize item}{\textcolor{OtagoGold}{\small$\blacktriangleright$}}
\setbeamertemplate{itemize subitem}{\textcolor{AccentBlue}{\small$\triangleright$}}
\setbeamertemplate{enumerate item}{\textcolor{OtagoBlue}{\small\bfseries\insertenumlabel.}}

\setbeamertemplate{title page}{%
	\begin{tikzpicture}[remember picture,overlay]
		\fill[OtagoBlue] (current page.north west)
		rectangle ([yshift=-0.38\paperheight]current page.north east);
	\end{tikzpicture}
	\vspace{0.06\paperheight}
	\begin{beamercolorbox}[wd=\textwidth, sep=0pt]{title}
		{\color{white}\Large\textbf{\inserttitle}}\\[0.4em]
		{\color{OtagoGold}\large\insertsubtitle}
	\end{beamercolorbox}
	\vspace{1.4em}
	\begin{beamercolorbox}[wd=\textwidth, sep=0pt]{author}
		{\color{TextGrey}\normalsize\insertauthor}\\[0.2em]
		{\color{TextGrey}\small\insertinstitute}\\[0.2em]
		{\color{TextGrey}\small\insertdate}
\end{beamercolorbox}}

% ─── Packages ──────────────────────────────────────────────
\usepackage{tikz}
\usepackage{booktabs}
\usepackage{array}
\usepackage{graphicx}
\usepackage{hyperref}
\usepackage{amsmath}
\usepackage{amssymb}
\usepackage{colortbl}
\usepackage{multirow}
\usepackage{xcolor}
\hypersetup{colorlinks=true, linkcolor=AccentBlue, urlcolor=AccentBlue,
	citecolor=AccentBlue, pdfauthor={Austin Henderson},
	pdftitle={ECON306 S2 2026}}

% ─── Images folder ─────────────────────────────────────────
%  Place all slide images in the 'images/' subfolder next to each deck.
%  Usage: \includegraphics[width=0.85\textwidth]{figure_name}
\graphicspath{{images/}}

% ─── Reusable macros ───────────────────────────────────────
% Recap slide at start of each lecture
\newcommand{\recapslide}[1]{%
	\begin{frame}{Recap: What Did We Cover Last Time?}
		\begin{block}{Key points from last lecture}
			\begin{itemize}#1\end{itemize}
		\end{block}
\end{frame}}

% Plan slide at start of each lecture
\newcommand{\planslide}[2]{%
	\begin{frame}{Today's Plan — #1}
		\begin{columns}
			\column{0.6\textwidth}
			\begin{itemize}#2\end{itemize}
			\column{0.38\textwidth}
			\begin{alertblock}{Reading}
				\footnotesize See Brightspace (Aoroa) for required readings before this lecture.
			\end{alertblock}
		\end{columns}
\end{frame}}

% Summary slide at end of each lecture
\newcommand{\summaryside}[1]{%
	\begin{frame}{Lecture Summary}
		\begin{exampleblock}{Key takeaways}
			\begin{itemize}#1\end{itemize}
		\end{exampleblock}
		\vspace{0.5em}
		{\footnotesize\textcolor{OtagoGold}{\textbf{Brightspace (Aoroa):}} slides, readings, and any announcements posted there.}
\end{frame}}

% NZ spotlight box
\newcommand{\nzbox}[2]{%
	\begin{alertblock}{\includegraphics[height=0.8em]{nz_flag_icon}\hspace{0.3em} NZ Spotlight: #1}
		#2
\end{alertblock}}

% Tutorial callout
\newcommand{\tutorialbox}[1]{%
	\begin{block}{\textbf{Tutorial This Week}}
		#1
\end{block}}

% ─── Section dividers ──────────────────────────────────────
\AtBeginSection[]{
	\begin{frame}[plain]
		\begin{tikzpicture}[remember picture,overlay]
			\fill[OtagoBlue!95] (current page.north west)
			rectangle (current page.south east);
			\node[anchor=center, text=white, font=\Large\bfseries,
			text width=0.9\paperwidth, align=center]
			at (current page.center) {\insertsectionhead};
			\node[anchor=south, text=OtagoGold, font=\normalsize,
			text width=0.9\paperwidth, align=center]
			at ([yshift=2.5cm]current page.south) {ECON306 — University of Otago};
		\end{tikzpicture}
\end{frame}}


\title[ECON306: W11D1 — Signalling Theory]{ECON306\\The Economics of Health and Education}
\subtitle{W11D1: Signalling Theory (Spence 1973); Separating Equilibria}
\author{Austin Henderson}
\institute{Department of Economics\\University of Otago}
\date{Wednesday 30 September 2026}

\begin{document}

\begin{frame}[plain]
  \titlepage
\end{frame}

\recapslide{
  \item Human capital theory: education raises productivity $\Rightarrow$ positive wage premium
  \item NZ graduate premium $\approx$60--80\%; falls for older students (shorter horizon)
  \item Ability bias and selection bias complicate empirical estimation of returns
}

\planslide{Signalling Theory: An Alternative Explanation}{
  \item Human capital vs.\ signalling: the key distinction
  \item Spence (1973) signalling model: setup and logic
  \item Separating equilibrium: how education signals ability without raising it
  \item Welfare implications of signalling equilibria
}

\begin{frame}{The Big Question: Why Does Education Raise Wages?}
  \begin{columns}
    \column{0.5\textwidth}
      \begin{block}{Human capital view}
        \small Education raises workers' \textbf{productivity} (skills, knowledge, problem-solving).
        Employers pay more because graduates genuinely produce more.\\[0.2em]
        \textit{Policy implication}: subsidise education --- it raises social welfare.
      \end{block}
    \column{0.47\textwidth}
      \begin{alertblock}{Signalling view}
        \small Education \textbf{does not} raise productivity. Instead, it serves as a \textbf{signal}
        of pre-existing ability to employers who cannot observe ability directly.\\[0.2em]
        \textit{Policy implication}: subsidising education may be wasteful --- it just
        shifts the goalposts on the signalling game.
      \end{alertblock}
  \end{columns}
  \vspace{0.2em}
  \begin{exampleblock}{Both views consistent with the same data}
    \small Graduates earn more than non-graduates. Both theories predict this. Distinguishing them
    empirically is one of the central challenges of labour economics.
  \end{exampleblock}
\end{frame}

\begin{frame}{Spence (1973): The Signalling Model}
  \begin{columns}
    \column{0.52\textwidth}
      \footnotesize \textbf{Michael Spence --- Nobel Prize 2001} (with Akerlof and Stiglitz)\\[0.2em]
      Education may \textbf{not} raise productivity at all --- it may serve purely as a
      \textbf{signal} of pre-existing ability to employers.
      \begin{itemize}
        \item \textbf{Information problem}: employers cannot observe worker ability before
              hiring (asymmetric information)
        \item \textbf{Costly signal}: education is costly but \emph{more costly} for
              \textbf{low-ability} workers (they find it harder)
        \item \textbf{High-ability} workers signal through education; low-ability do not
        \item Result: a \textbf{separating equilibrium}
      \end{itemize}
    \column{0.44\textwidth}
      \begin{alertblock}{The crucial implication}
        \footnotesize If signalling theory is correct:
        \begin{itemize}
          \item Private returns are \textbf{positive} (wage premium is real)
          \item But \textbf{social returns are zero} (education just sorts workers)
          \item Government subsidy would be a \textbf{waste of resources}
        \end{itemize}
      \end{alertblock}
  \end{columns}
\end{frame}

\begin{frame}{The Separating Equilibrium: Setup}
  \begin{block}{Setup (simplified)}
    \small Two types of workers: \textbf{High ability} ($H$) and \textbf{Low ability} ($L$).
    \begin{itemize}
      \item $H$ workers produce \$150K/year; $L$ workers produce \$80K/year
      \item Employers cannot tell them apart before hiring
      \item Education costs \textbf{\$10K/year} for $H$ but \textbf{\$25K/year} for $L$
            (low-ability workers find education harder/costlier)
    \end{itemize}
  \end{block}
\end{frame}

\begin{frame}{The Separating Equilibrium: The Logic}
  \begin{block}{The separating equilibrium}
    \small Suppose employers pay \$150K to graduates, \$80K to non-graduates:
    \begin{itemize}
      \item $H$ workers: benefit $=$ \$70K; cost of 4 years $= 4\times\$10\text{K} = \$40\text{K}$.
            \textbf{Net gain $=$ \$30K. Get the degree.}
      \item $L$ workers: benefit $=$ \$70K; cost $= 4\times\$25\text{K} = \$100\text{K}$.
            \textbf{Net loss $=$ -\$30K. Don't get the degree.}
      \item Employers' beliefs are confirmed. \textbf{Equilibrium.}
    \end{itemize}
  \end{block}
  \vspace{0.3em}
  \small \textbf{Key insight}: Education doesn't need to teach anything to sustain this equilibrium
  --- it only needs to be differentially costly across worker types.
\end{frame}

\begin{frame}{Key Properties of the Signalling Equilibrium}
  \small
  \begin{enumerate}
    \item \textbf{Stability}: no agent wants to deviate given others' beliefs
    \begin{itemize}
      \item $H$ workers don't want to stop educating (they'd be paid \$80K)
      \item $L$ workers don't want to start educating (cost exceeds benefit)
      \item Employers' wage-setting is confirmed by actual worker quality
    \end{itemize}
    \item \textbf{Pooling equilibrium alternative}: if employers cannot distinguish, they pay
          the average wage; but then $H$ workers have incentive to signal
    \item \textbf{Welfare implications}:
    \begin{itemize}
      \item Total wages paid = same as without signalling (productivity unchanged)
      \item But resources are wasted on education that produced no skills
      \item Real education consumes resources; sorting is costly
    \end{itemize}
    \item \textbf{``Arms race'' dynamic}: if everyone gets more education, the signal
          threshold may rise, leaving everyone worse off (signalling inflation)
  \end{enumerate}
\end{frame}

\begin{frame}{Conditions for a Signalling Equilibrium}
  \begin{columns}
    \column{0.52\textwidth}
      \begin{alertblock}{The single-crossing property}
        \small For a separating equilibrium to exist:
        \[\text{Cost}_L > \text{Wage gap} > \text{Cost}_H\]
        Indifference curves of $H$ and $L$ workers cross only once in education--wage space.
      \end{alertblock}
    \column{0.45\textwidth}
      \begin{block}{What makes education a good signal?}
        \small
        \begin{itemize}
          \item Costly to fake (time and effort even if it teaches little)
          \item High-ability workers find it less costly (stronger performance, faster completion)
          \item Credentials hard to counterfeit (graded exams, external assessment)
        \end{itemize}
      \end{block}
  \end{columns}
\end{frame}

\begin{frame}{Policy Implications of Signalling Theory}
  \begin{columns}
    \column{0.52\textwidth}
      \begin{alertblock}{If signalling theory is mostly true:}
        \footnotesize
        \begin{itemize}
          \item Subsidising tertiary education accelerates signalling inflation
          \item Private returns positive; social returns $\approx 0$
          \item Government subsidy is welfare-reducing (resources wasted)
          \item \textbf{Policy}: reduce subsidies; let markets sort
        \end{itemize}
      \end{alertblock}
    \column{0.45\textwidth}
      \begin{block}{If human capital theory is mostly true:}
        \footnotesize
        \begin{itemize}
          \item Education genuinely raises productivity and social welfare
          \item Externalities justify subsidy beyond capital market correction
          \item \textbf{Policy}: subsidise, especially for low-income students
        \end{itemize}
      \end{block}
  \end{columns}
  \vspace{0.15em}
  \begin{exampleblock}{The honest answer}
    \small Both theories are partially true. The NZ policy debate is about the \emph{relative}
    magnitudes --- examined in W11D2 and W11D3.
  \end{exampleblock}
\end{frame}

\summaryside{
  \footnotesize
  \item Spence (1973) signalling model: education may signal pre-existing ability without raising productivity; Nobel Prize 2001
  \item Separating equilibrium: high-ability workers signal with education; low-ability do not, because signalling costs differ
  \item Single-crossing property: cost of education must be lower for high-ability workers than the wage gap to sustain equilibrium
  \item Policy implication: if pure signalling, subsidising education wastes resources and inflates the signalling threshold
}

\end{document}
